\begin{table}[!htbp]
\centering\small
\caption{Związek miar ze średnimi rangami wariantów: korelacja Spearmana między średnią wartością miary wariantu (po zbiorach) a~jego średnią rangą dla 15 wariantów. Korelacja ujemna oznacza, że wariant o~wyższej wartości miary ma lepszą rangę. Wartości $p_H$ podano po korekcie Holma.}
\label{tab:miary-rangi}
\begin{tabular}{lrrrr}
\toprule
& \multicolumn{2}{c}{\textbf{Klasyfikacja}} & \multicolumn{2}{c}{\textbf{Regresja}} \\
\cmidrule(lr){2-3}\cmidrule(lr){4-5}
\textbf{Miara} & $\rho$ & $p_H$ & $\rho$ & $p_H$ \\
\midrule
$D_J$ & +0.55 & 0.067 & +0.17 & 1.000 \\
$D_C$ & +0.45 & 0.089 & -0.01 & 1.000 \\
$D_Y$ & +0.64 & \textbf{0.031} & +0.21 & 1.000 \\
$U$ & -0.90 & \textbf{<0.001} & -0.70 & \textbf{0.027} \\
$R_Y$ & +0.79 & \textbf{0.003} & +0.42 & 0.592 \\
$R_X$ & +0.78 & \textbf{0.003} & +0.55 & 0.203 \\
różnorodność predykcji & +0.66 & \textbf{0.031} & +0.11 & 1.000 \\
\bottomrule
\end{tabular}
\end{table}
